%%%%%%%%%%%%%%%%%%%%%%%%%%%%%%%%%%%%%%%%%%%%%%%%%%%%%%%%%%%%%%%%%%%%%%
% Lớp 10 | Xác suất | Xác suất của biến cố
% Dạng 0D0?2-1: Các câu hỏi lý thuyết tổng hợp
% Số câu: 8  (sinh tự động bởi tach_cau.py; từ nay thầy sửa trực tiếp file này)
%%%%%%%%%%%%%%%%%%%%%%%%%%%%%%%%%%%%%%%%%%%%%%%%%%%%%%%%%%%%%%%%%%%%%%

%%%%%%%%%%%%%%%% TRẮC NGHIỆM 4 PHƯƠNG ÁN %%%%%%%%%%%%%%%%

% nguồn: MaHoaNH_TracNghiem_Lop10
\begin{ex}%[0D0N2-1]%[ma:3e77c19e]
	Xét một phép thử có không gian mẫu $\Omega$ và $A$ là một biến cố của phép thử đó.  Phát biểu nào dưới đây là \textbf{sai}?
	\choice 
		{\True $\mathrm{P}(A)=0$ khi và  chỉ khi $A$ là chắc chắn}
		{$0 \leq \mathrm{P}(A)  \leq 1$}
		{$\mathrm{P}(A)=1-\mathrm{P}\left(  \overline{A}\right) $}
		{$\mathrm{P}(A)=\dfrac{\mathrm{n} (A)}{\mathrm{n} (\Omega)}$} 
	 \loigiai{
		Ta có $\mathrm{P}(A)=0$ khi và chỉ khi $A$ là biến cố không thể xảy ra nên khẳng định sai là  $\mathrm{P}(A)=0$ khi và chỉ khi $A$ là chắc chắn.
	} 
\end{ex}

% nguồn: MaHoaNH_TracNghiem_Lop10
\begin{ex}%[0D0N2-1]%[ma:700038d0]
	Cho  phép thử $T$ có không gian mẫu là $\Omega$. Gọi $A$ là một biến cố liên quan đến phép thử $T$. Giả sử $n(A) = 4$; $n(\Omega) = 13$ thì xác suất của biến cố $A$ bằng
	\choice 
		{$\dfrac{13}{4}$} 
		{\True $\dfrac{4}{13}$} 
		{$4\cdot  13$}
		 {$13-4$}
	 \loigiai{
		Vì $n(A)  = 4$; $n(\Omega) = 13$ nên xác suất của biến cố $A$ bằng $\mathrm{P}(A) = \dfrac{n(A)}{n(\Omega)} = \dfrac{4}{13}$.
	} 
\end{ex}

% nguồn: MaHoaNH_TracNghiem_Lop10
\begin{ex}%[0D0H2-1]%[ma:4f315780]
	Một em bé có bộ $6$ thẻ chữ, trên mỗi thẻ có ghi một chữ cái, trong đó có $3$ thẻ chữ T, một thẻ chữ N, một thẻ chữ H và một thẻ chữ P. Em bé đó xếp ngẫu nhiên $6$ thẻ đó thành một hàng ngang. Tính xác suất em bé xếp được thành  dãy \textbf{TNTHPT}.
	 \choice
		{\True $\dfrac{1}{120}$} 
		{$\dfrac{1}{720}$} 
		{$\dfrac{1}{6}$} 
		{$\dfrac{1}{20}$} 
	 \loigiai{
		Xem ba chữ T là phân biệt ta  có $n(\Omega) = 6!$.\\
		Gọi $A$ là biến cố \lq\lq xếp ngẫu nhiên $6$ thẻ đó thành dãy  \textbf{TNTHPT} \rq\rq.\\
		Suy ra  $n(A) = 3!$ (vì hoán vị các chữ T, còn N, H, P cố định).\\
		Vậy xác suất biến cố  $A$ là $P(A) = \dfrac{3!}{6!} = \dfrac{1}{120}$.
	} 
\end{ex}

% nguồn: MaHoaNH_TracNghiem_Lop10
\begin{ex}%[0D0H2-1]%[ma:5fba5733]
	Một bộ bài tú lơ khơ gồm $52$ quân. Lấy ngẫu nhiên $2$ quân bài. Xác suất lấy được $2$ quân át  bằng
	 \choice
		 {\True $\dfrac{1}{121}$}
		 {$\dfrac{2}{663}$}
		{$\dfrac{1}{1326}$} 
		 {$\dfrac{1}{26}$}
	 \loigiai{
		Số phần tử  của không gian mẫu là số cách chọn $2$ lá từ $52$ lá
		\[ 
		n(\Omega) =  \mathrm{C}^2_{52}.
		\] 
		Biến cố $A$: \lq\lq Lấy được $2$ quân át\rq\rq $\Rightarrow n(A) =  \mathrm{C}^2_4$.\\
		Vậy xác suất  là
		\[ 
		 \mathrm{P}(A) = \dfrac{\mathrm{C}^2_4}{\mathrm{C}^2_{52}} = \dfrac{6}{1326} = \dfrac{1}{121}.
		\] 
	 }
\end{ex}

% nguồn: MaHoaNH_TracNghiem_Lop10
\begin{ex}%[0D0V2-1]%[ma:5b89e8fd]
	Gieo một con xúc xắc cân đối và đồng chất. Giả sử xúc xắc xuất hiện mặt $b$ chấm. Tính xác suất để phương trình  $x^2 + 2bx + 4 = 0$ có nghiệm.
	\choice 
		{\True  $\dfrac{5}{6}$}
		{$\dfrac{1}{7}$} 
		 {$\dfrac{2}{9}$}
		{$\dfrac{3}{8}$} 
	\loigiai{ 
		Theo giả thiết ta có  $b \in \{1;2;3;4;5;6\}$.\\
		Tập hợp $\Omega$ có $6$  phần tử.\\
		Gọi $A$ là biến cố \lq\lq gieo xúc xắc để phương trình $x^2 + 2bx  + 4 = 0$ có nghiệm\rq\rq.\\
		Phương trình có nghiệm khi $\Delta = 4b^2 - 16 \ge 0 \Leftrightarrow b^2 \ge 4 \Leftrightarrow  b \ge 2$.\\
		Do  đó $b \in \{2;3;4;5;6\} \Rightarrow n(A) = 5$.\\
		Vậy xác suất là $P(A) =  \dfrac{n(A)}{n(\Omega)} = \dfrac{5}{6}$.
	 }
\end{ex}


%%%%%%%%%%%%%%%% ĐÚNG/SAI %%%%%%%%%%%%%%%%

% nguồn: MaHoaNH_DungSai_Lop10
\begin{ex}%[0D0N2-1]%[ma:1e59897f]
	Khi thực hiện phép thử $T$ chỉ có một số hữu hạn các kết quả đồng khả năng xuất hiện. Gọi $n(\Omega)$ là số kết quả có thể xảy ra của phép thử, $A$ là biến cố liên quan đến phép thử $T$, $n(A)$ là số kết quả thuận lợi cho biến cố $A$, $P(A)$ là  xác suất của biến cố $A$. Các khẳng định sau đúng hay sai?
	\choiceTF 
		{$P(A)=n(\Omega)$} 
		{$P(A)=\dfrac{n(\Omega)}{n(A)}$} 
		{\True $ 0\leq \mathrm{P}(A)\leq 1  $}
		{\True  $P(A)=\dfrac{n(A)}{n(\Omega)}$}
	 \loigiai{
		 Theo định nghĩa xác suất $P(A)=\dfrac{n(A)}{n(\Omega)}$.\\
		$ 0\leq \mathrm{P}(A)\leq 1$. 
	} 
\end{ex}

% nguồn: MaHoaNH_DungSai_Lop10
\begin{ex}%[0D0N2-1]%[ma:77fcfc1a]
	Giả sử $ A $ là biến cố liên quan đến một  phép thử có không gian mẫu là $ \Omega $. Các khẳng định sau đúng hay sai?
	 \choiceTF
		{\True $  \mathrm{P}(\Omega)=1 $}
		 {$\mathrm{P}(A)=\mathrm{P}(A)-1 $}
		{\True $ 0\leq \mathrm{P}(A)\leq 1 $} 
		{$\mathrm{P}=n(\varnothing)$} 
	 \loigiai{
		$\mathrm{P}(\Omega)=1$ là khẳng định đúng.\\ 
		$ 0\leq \mathrm{P}(A)\leq 1 $ là khẳng định đúng.\\ 
		Ta có $ \mathrm{P}\left(\overline{A}\right)=1-\mathrm{P}(A)  $. Suy ra
		$ \mathrm{P}\left(\overline{A}\right)=\mathrm{P}(A)-1 $ là khẳng  định sai.\\
		$\mathrm{P}(A)=n(\varnothing)$  là khẳng định sai.
	 }
\end{ex}

% nguồn: MaHoaNH_DungSai_Lop10
\begin{ex}%[0D0H2-1]%[ma:4d058834]
	Cho $A$ là một biến cố liên quan phép thử $T$. Chọn  đúng, sai cho mỗi khẳng định sau:
	 \choiceTF
		{$\mathrm{P}(A)>0$} 
		{\True  $\mathrm{P}(A) = 1 - \mathrm{P}\left(\overline{A}\right)$}
		{\True $\mathrm{P}(A) = 1 \Leftrightarrow A = \Omega$} 
		{$\mathrm{P}(A)<1$} 
	 \loigiai{
		 Với $A$ là một biến cố liên quan phép thử $T$. Khi đó
		\begin{itemize}
			\item $0 \leq  \mathrm{P}(A) \leq 1$ nên mệnh đề $\mathrm{P}(A)$ là số lớn hơn $0$ và $\mathrm{P}(A)$ là số nhỏ hơn $1$ là các mệnh đề sai.
			\item $\mathrm{P}(A) = 0  \Rightarrow A = \Omega$ nên mệnh đề $\mathrm{P}(A) = 0 \Leftrightarrow A = \Omega$ đúng.
			\item Mệnh đề $\mathrm{P}(A) = 1 -  \mathrm{P}(\overline{A})$ là mệnh đề đúng.
		\end{itemize}
	} 
\end{ex}
